% Options for packages loaded elsewhere
\PassOptionsToPackage{unicode}{hyperref}
\PassOptionsToPackage{hyphens}{url}
\PassOptionsToPackage{dvipsnames,svgnames,x11names}{xcolor}
%
\documentclass[
  letterpaper,
  DIV=11,
  numbers=noendperiod]{scrartcl}

\usepackage{amsmath,amssymb}
\usepackage{iftex}
\ifPDFTeX
  \usepackage[T1]{fontenc}
  \usepackage[utf8]{inputenc}
  \usepackage{textcomp} % provide euro and other symbols
\else % if luatex or xetex
  \usepackage{unicode-math}
  \defaultfontfeatures{Scale=MatchLowercase}
  \defaultfontfeatures[\rmfamily]{Ligatures=TeX,Scale=1}
\fi
\usepackage{lmodern}
\ifPDFTeX\else  
    % xetex/luatex font selection
\fi
% Use upquote if available, for straight quotes in verbatim environments
\IfFileExists{upquote.sty}{\usepackage{upquote}}{}
\IfFileExists{microtype.sty}{% use microtype if available
  \usepackage[]{microtype}
  \UseMicrotypeSet[protrusion]{basicmath} % disable protrusion for tt fonts
}{}
\makeatletter
\@ifundefined{KOMAClassName}{% if non-KOMA class
  \IfFileExists{parskip.sty}{%
    \usepackage{parskip}
  }{% else
    \setlength{\parindent}{0pt}
    \setlength{\parskip}{6pt plus 2pt minus 1pt}}
}{% if KOMA class
  \KOMAoptions{parskip=half}}
\makeatother
\usepackage{xcolor}
\setlength{\emergencystretch}{3em} % prevent overfull lines
\setcounter{secnumdepth}{-\maxdimen} % remove section numbering
% Make \paragraph and \subparagraph free-standing
\makeatletter
\ifx\paragraph\undefined\else
  \let\oldparagraph\paragraph
  \renewcommand{\paragraph}{
    \@ifstar
      \xxxParagraphStar
      \xxxParagraphNoStar
  }
  \newcommand{\xxxParagraphStar}[1]{\oldparagraph*{#1}\mbox{}}
  \newcommand{\xxxParagraphNoStar}[1]{\oldparagraph{#1}\mbox{}}
\fi
\ifx\subparagraph\undefined\else
  \let\oldsubparagraph\subparagraph
  \renewcommand{\subparagraph}{
    \@ifstar
      \xxxSubParagraphStar
      \xxxSubParagraphNoStar
  }
  \newcommand{\xxxSubParagraphStar}[1]{\oldsubparagraph*{#1}\mbox{}}
  \newcommand{\xxxSubParagraphNoStar}[1]{\oldsubparagraph{#1}\mbox{}}
\fi
\makeatother

\usepackage{color}
\usepackage{fancyvrb}
\newcommand{\VerbBar}{|}
\newcommand{\VERB}{\Verb[commandchars=\\\{\}]}
\DefineVerbatimEnvironment{Highlighting}{Verbatim}{commandchars=\\\{\}}
% Add ',fontsize=\small' for more characters per line
\usepackage{framed}
\definecolor{shadecolor}{RGB}{241,243,245}
\newenvironment{Shaded}{\begin{snugshade}}{\end{snugshade}}
\newcommand{\AlertTok}[1]{\textcolor[rgb]{0.68,0.00,0.00}{#1}}
\newcommand{\AnnotationTok}[1]{\textcolor[rgb]{0.37,0.37,0.37}{#1}}
\newcommand{\AttributeTok}[1]{\textcolor[rgb]{0.40,0.45,0.13}{#1}}
\newcommand{\BaseNTok}[1]{\textcolor[rgb]{0.68,0.00,0.00}{#1}}
\newcommand{\BuiltInTok}[1]{\textcolor[rgb]{0.00,0.23,0.31}{#1}}
\newcommand{\CharTok}[1]{\textcolor[rgb]{0.13,0.47,0.30}{#1}}
\newcommand{\CommentTok}[1]{\textcolor[rgb]{0.37,0.37,0.37}{#1}}
\newcommand{\CommentVarTok}[1]{\textcolor[rgb]{0.37,0.37,0.37}{\textit{#1}}}
\newcommand{\ConstantTok}[1]{\textcolor[rgb]{0.56,0.35,0.01}{#1}}
\newcommand{\ControlFlowTok}[1]{\textcolor[rgb]{0.00,0.23,0.31}{\textbf{#1}}}
\newcommand{\DataTypeTok}[1]{\textcolor[rgb]{0.68,0.00,0.00}{#1}}
\newcommand{\DecValTok}[1]{\textcolor[rgb]{0.68,0.00,0.00}{#1}}
\newcommand{\DocumentationTok}[1]{\textcolor[rgb]{0.37,0.37,0.37}{\textit{#1}}}
\newcommand{\ErrorTok}[1]{\textcolor[rgb]{0.68,0.00,0.00}{#1}}
\newcommand{\ExtensionTok}[1]{\textcolor[rgb]{0.00,0.23,0.31}{#1}}
\newcommand{\FloatTok}[1]{\textcolor[rgb]{0.68,0.00,0.00}{#1}}
\newcommand{\FunctionTok}[1]{\textcolor[rgb]{0.28,0.35,0.67}{#1}}
\newcommand{\ImportTok}[1]{\textcolor[rgb]{0.00,0.46,0.62}{#1}}
\newcommand{\InformationTok}[1]{\textcolor[rgb]{0.37,0.37,0.37}{#1}}
\newcommand{\KeywordTok}[1]{\textcolor[rgb]{0.00,0.23,0.31}{\textbf{#1}}}
\newcommand{\NormalTok}[1]{\textcolor[rgb]{0.00,0.23,0.31}{#1}}
\newcommand{\OperatorTok}[1]{\textcolor[rgb]{0.37,0.37,0.37}{#1}}
\newcommand{\OtherTok}[1]{\textcolor[rgb]{0.00,0.23,0.31}{#1}}
\newcommand{\PreprocessorTok}[1]{\textcolor[rgb]{0.68,0.00,0.00}{#1}}
\newcommand{\RegionMarkerTok}[1]{\textcolor[rgb]{0.00,0.23,0.31}{#1}}
\newcommand{\SpecialCharTok}[1]{\textcolor[rgb]{0.37,0.37,0.37}{#1}}
\newcommand{\SpecialStringTok}[1]{\textcolor[rgb]{0.13,0.47,0.30}{#1}}
\newcommand{\StringTok}[1]{\textcolor[rgb]{0.13,0.47,0.30}{#1}}
\newcommand{\VariableTok}[1]{\textcolor[rgb]{0.07,0.07,0.07}{#1}}
\newcommand{\VerbatimStringTok}[1]{\textcolor[rgb]{0.13,0.47,0.30}{#1}}
\newcommand{\WarningTok}[1]{\textcolor[rgb]{0.37,0.37,0.37}{\textit{#1}}}

\providecommand{\tightlist}{%
  \setlength{\itemsep}{0pt}\setlength{\parskip}{0pt}}\usepackage{longtable,booktabs,array}
\usepackage{calc} % for calculating minipage widths
% Correct order of tables after \paragraph or \subparagraph
\usepackage{etoolbox}
\makeatletter
\patchcmd\longtable{\par}{\if@noskipsec\mbox{}\fi\par}{}{}
\makeatother
% Allow footnotes in longtable head/foot
\IfFileExists{footnotehyper.sty}{\usepackage{footnotehyper}}{\usepackage{footnote}}
\makesavenoteenv{longtable}
\usepackage{graphicx}
\makeatletter
\def\maxwidth{\ifdim\Gin@nat@width>\linewidth\linewidth\else\Gin@nat@width\fi}
\def\maxheight{\ifdim\Gin@nat@height>\textheight\textheight\else\Gin@nat@height\fi}
\makeatother
% Scale images if necessary, so that they will not overflow the page
% margins by default, and it is still possible to overwrite the defaults
% using explicit options in \includegraphics[width, height, ...]{}
\setkeys{Gin}{width=\maxwidth,height=\maxheight,keepaspectratio}
% Set default figure placement to htbp
\makeatletter
\def\fps@figure{htbp}
\makeatother

\KOMAoption{captions}{tableheading}
\makeatletter
\@ifpackageloaded{caption}{}{\usepackage{caption}}
\AtBeginDocument{%
\ifdefined\contentsname
  \renewcommand*\contentsname{Table of contents}
\else
  \newcommand\contentsname{Table of contents}
\fi
\ifdefined\listfigurename
  \renewcommand*\listfigurename{List of Figures}
\else
  \newcommand\listfigurename{List of Figures}
\fi
\ifdefined\listtablename
  \renewcommand*\listtablename{List of Tables}
\else
  \newcommand\listtablename{List of Tables}
\fi
\ifdefined\figurename
  \renewcommand*\figurename{Figure}
\else
  \newcommand\figurename{Figure}
\fi
\ifdefined\tablename
  \renewcommand*\tablename{Table}
\else
  \newcommand\tablename{Table}
\fi
}
\@ifpackageloaded{float}{}{\usepackage{float}}
\floatstyle{ruled}
\@ifundefined{c@chapter}{\newfloat{codelisting}{h}{lop}}{\newfloat{codelisting}{h}{lop}[chapter]}
\floatname{codelisting}{Listing}
\newcommand*\listoflistings{\listof{codelisting}{List of Listings}}
\makeatother
\makeatletter
\makeatother
\makeatletter
\@ifpackageloaded{caption}{}{\usepackage{caption}}
\@ifpackageloaded{subcaption}{}{\usepackage{subcaption}}
\makeatother

\ifLuaTeX
  \usepackage{selnolig}  % disable illegal ligatures
\fi
\usepackage{bookmark}

\IfFileExists{xurl.sty}{\usepackage{xurl}}{} % add URL line breaks if available
\urlstyle{same} % disable monospaced font for URLs
\hypersetup{
  pdftitle={Data Structures},
  pdfauthor={Joseph Mhango},
  colorlinks=true,
  linkcolor={blue},
  filecolor={Maroon},
  citecolor={Blue},
  urlcolor={Blue},
  pdfcreator={LaTeX via pandoc}}


\title{Data Structures}
\usepackage{etoolbox}
\makeatletter
\providecommand{\subtitle}[1]{% add subtitle to \maketitle
  \apptocmd{\@title}{\par {\large #1 \par}}{}{}
}
\makeatother
\subtitle{Data structures}
\author{Joseph Mhango}
\date{2024-09-13}

\begin{document}
\maketitle


\subsection{Objectives}\label{objectives}

By the end of this section, you should be able to:

\begin{itemize}
\tightlist
\item
  Identify and describe basic data types in R.
\item
  Understand how factors represent categorical data.
\item
  Use \texttt{class()} and conversion functions to change variable
  types.
\item
  Work with vectors, matrices, and lists.
\item
  Apply best practices in naming and managing R objects.
\end{itemize}

\begin{center}\rule{0.5\linewidth}{0.5pt}\end{center}

\subsection{Basic data types in R}\label{basic-data-types-in-r}

R works with different kinds of \textbf{data objects}, and each object
has an associated \textbf{type}. These types determine what operations
are possible and how the data is treated during analysis. All the data
you create lives in the \textbf{Global Environment}, and it's important
to know what kind of data you're working with.

R is pretty good at guessing what type your data should be --- but this
can sometimes backfire. Think of R as a \emph{passive-aggressive
butler}: it tries to help, but doesn't always tell you when it's done
something ``helpful.'' So, you should learn to check and control data
types explicitly.

\begin{center}\rule{0.5\linewidth}{0.5pt}\end{center}

\subsection{Basic primitive types in
R}\label{basic-primitive-types-in-r}

Here are the most common \textbf{primitive data types}:

\begin{Shaded}
\begin{Highlighting}[]
\NormalTok{a }\OtherTok{\textless{}{-}} \DecValTok{1}    \CommentTok{\# numeric}
\FunctionTok{mode}\NormalTok{(a)}
\end{Highlighting}
\end{Shaded}

\begin{verbatim}
[1] "numeric"
\end{verbatim}

\begin{Shaded}
\begin{Highlighting}[]
\NormalTok{a }\OtherTok{\textless{}{-}} \StringTok{"1"}  \CommentTok{\# character}
\FunctionTok{mode}\NormalTok{(a)}
\end{Highlighting}
\end{Shaded}

\begin{verbatim}
[1] "character"
\end{verbatim}

\begin{Shaded}
\begin{Highlighting}[]
\NormalTok{a }\OtherTok{\textless{}{-}} \StringTok{\textquotesingle{}1\textquotesingle{}}  \CommentTok{\# character}
\FunctionTok{mode}\NormalTok{(a)}
\end{Highlighting}
\end{Shaded}

\begin{verbatim}
[1] "character"
\end{verbatim}

\begin{Shaded}
\begin{Highlighting}[]
\NormalTok{a }\OtherTok{\textless{}{-}} \ConstantTok{TRUE} \CommentTok{\# logical}
\FunctionTok{mode}\NormalTok{(a)}
\end{Highlighting}
\end{Shaded}

\begin{verbatim}
[1] "logical"
\end{verbatim}

\begin{Shaded}
\begin{Highlighting}[]
\NormalTok{a }\OtherTok{\textless{}{-}}  \DecValTok{1}\NormalTok{i  }\CommentTok{\# complex}
\FunctionTok{mode}\NormalTok{(a)}
\end{Highlighting}
\end{Shaded}

\begin{verbatim}
[1] "complex"
\end{verbatim}

These types include: - \textbf{numeric}: Numbers like 1, 3.14 -
\textbf{character}: Strings like ``apple'' - \textbf{logical}: TRUE or
FALSE - \textbf{complex}: Numbers with imaginary parts like 1i

\begin{center}\rule{0.5\linewidth}{0.5pt}\end{center}

\subsection{Special values in R}\label{special-values-in-r}

There are a few special types of values that you'll come across often:

\begin{itemize}
\tightlist
\item
  \texttt{NULL} means an undefined or empty object (a missing vector).
  Use \texttt{is.null()} to test it.
\item
  \texttt{NA} stands for \emph{Not Available} and represents missing
  data. Check with \texttt{is.na()}.
\item
  \texttt{Inf} and \texttt{-Inf} mean positive or negative infinity
  (e.g., from \texttt{1/0}). Use \texttt{is.infinite()}.
\item
  \texttt{NaN} means \emph{Not a Number}, e.g.~\texttt{0/0}. Use
  \texttt{is.nan()}.
\end{itemize}

\begin{Shaded}
\begin{Highlighting}[]
\NormalTok{my\_list }\OtherTok{\textless{}{-}} \FunctionTok{list}\NormalTok{(}\AttributeTok{a =} \DecValTok{1}\NormalTok{, }\AttributeTok{b =} \DecValTok{2}\NormalTok{)}
\NormalTok{my\_list }\OtherTok{\textless{}{-}}\NormalTok{ my\_list}\SpecialCharTok{$}\NormalTok{c  }\CommentTok{\# Element \textquotesingle{}c\textquotesingle{} doesn\textquotesingle{}t exist, Output: NULL}

\NormalTok{na\_addition  }\OtherTok{\textless{}{-}} \DecValTok{5} \SpecialCharTok{+} \ConstantTok{NA}
\FunctionTok{print}\NormalTok{(na\_addition)  }\CommentTok{\# Output: NA}
\end{Highlighting}
\end{Shaded}

\begin{verbatim}
[1] NA
\end{verbatim}

\begin{center}\rule{0.5\linewidth}{0.5pt}\end{center}

\subsection{Naming objects in R}\label{naming-objects-in-r}

There are some general rules and conventions for naming objects in R:

\begin{itemize}
\tightlist
\item
  Names must start with a letter.
\item
  Names can include letters, numbers, underscores (\texttt{\_}), and
  dots (\texttt{.}).
\item
  Avoid special characters like \texttt{@}, \texttt{\%}, \texttt{\#},
  etc.
\item
  Avoid spaces --- use underscores or camelCase instead.
\item
  Use names that are short, descriptive, and consistent.
\end{itemize}

\begin{Shaded}
\begin{Highlighting}[]
\NormalTok{x1 }\OtherTok{\textless{}{-}} \DecValTok{5}  

\NormalTok{x2 }\OtherTok{\textless{}{-}} \StringTok{"It was a dark and stormy night"} 

\CommentTok{\# harder to read and maintain:}
\NormalTok{my\_variable\_9283467 }\OtherTok{\textless{}{-}} \DecValTok{1} 
\end{Highlighting}
\end{Shaded}

\begin{center}\rule{0.5\linewidth}{0.5pt}\end{center}

\subsection{Data with factors}\label{data-with-factors}

Factors are how R represents \textbf{categorical data}. This is useful
when you have data that falls into named categories, such as treatment
groups, soil types, or crop varieties.

Let's say you want to analyse categorical data. Sometimes this is called
a \emph{factor}.

\includegraphics{pics/mean_mark_boxplot_ggplot2.webp}

There are two types of factors:

\begin{itemize}
\tightlist
\item
  \textbf{Unordered (nominal)}: No implied order among categories.

  \begin{itemize}
  \tightlist
  \item
    Example:
    \texttt{cow\_names\ \textless{}-\ factor(c("Martha",\ "Ruth",\ "Edith"))}
  \end{itemize}
\item
  \textbf{Ordered}: There is a logical order.

  \begin{itemize}
  \tightlist
  \item
    Example:
    \texttt{dose\ \textless{}-\ factor(c("Control",\ "Half",\ "Full"),\ ordered\ =\ TRUE)}
  \end{itemize}
\end{itemize}

\begin{center}\rule{0.5\linewidth}{0.5pt}\end{center}

\subsection{\texorpdfstring{Using \texttt{class()} and converting
variables}{Using class() and converting variables}}\label{using-class-and-converting-variables}

You can check the type of an object with \texttt{class()} or
\texttt{mode()}, and convert between types using \texttt{as.*}
functions.

\begin{Shaded}
\begin{Highlighting}[]
\CommentTok{\# non{-}ordered factor}
\NormalTok{variety }\OtherTok{\textless{}{-}} \FunctionTok{c}\NormalTok{(}\StringTok{"short"}\NormalTok{, }\StringTok{"short"}\NormalTok{, }\StringTok{"short"}\NormalTok{,}
             \StringTok{"early"}\NormalTok{, }\StringTok{"early"}\NormalTok{, }\StringTok{"early"}\NormalTok{,}
             \StringTok{"hybrid"}\NormalTok{, }\StringTok{"hybrid"}\NormalTok{, }\StringTok{"hybrid"}\NormalTok{)}

\NormalTok{factor\_variety }\OtherTok{\textless{}{-}} \FunctionTok{factor}\NormalTok{(variety)  }\CommentTok{\# creates an unordered factor}
\FunctionTok{class}\NormalTok{(factor\_variety)}
\end{Highlighting}
\end{Shaded}

\begin{verbatim}
[1] "factor"
\end{verbatim}

\begin{center}\rule{0.5\linewidth}{0.5pt}\end{center}

\subsection{Vector and Matrix fun}\label{vector-and-matrix-fun}

These are the fundamental data containers in R. They organise values in
structured ways:

\begin{itemize}
\tightlist
\item
  \textbf{Vector}: One-dimensional, holds values of the same type.
  Indexed as \texttt{my\_vec{[}i{]}}
\item
  \textbf{Matrix}: Two-dimensional, also holds only one type. Indexed as
  \texttt{my\_mat{[}i,\ j{]}}
\item
  \textbf{Array}: Extends matrix to 3+ dimensions, indexed as
  \texttt{my\_array{[}i,\ j,\ k{]}}
\end{itemize}

\begin{Shaded}
\begin{Highlighting}[]
\CommentTok{\# Vectors}
\NormalTok{myvec1 }\OtherTok{\textless{}{-}} \FunctionTok{c}\NormalTok{(}\DecValTok{1}\NormalTok{,}\DecValTok{2}\NormalTok{,}\DecValTok{3}\NormalTok{,}\DecValTok{4}\NormalTok{,}\DecValTok{5}\NormalTok{) }
\NormalTok{myvec1}
\end{Highlighting}
\end{Shaded}

\begin{verbatim}
[1] 1 2 3 4 5
\end{verbatim}

\begin{Shaded}
\begin{Highlighting}[]
\FunctionTok{class}\NormalTok{(myvec1) }
\end{Highlighting}
\end{Shaded}

\begin{verbatim}
[1] "numeric"
\end{verbatim}

\begin{Shaded}
\begin{Highlighting}[]
\NormalTok{myvec2 }\OtherTok{\textless{}{-}} \FunctionTok{as.character}\NormalTok{(myvec1) }
\end{Highlighting}
\end{Shaded}

\begin{Shaded}
\begin{Highlighting}[]
\CommentTok{\# Matrices}
\NormalTok{vec1 }\OtherTok{\textless{}{-}} \DecValTok{1}\SpecialCharTok{:}\DecValTok{16} 
\NormalTok{vec1 }
\end{Highlighting}
\end{Shaded}

\begin{verbatim}
 [1]  1  2  3  4  5  6  7  8  9 10 11 12 13 14 15 16
\end{verbatim}

\begin{Shaded}
\begin{Highlighting}[]
\NormalTok{mat1 }\OtherTok{\textless{}{-}} \FunctionTok{matrix}\NormalTok{(}\AttributeTok{data =}\NormalTok{ vec1, }
               \AttributeTok{ncol =} \DecValTok{4}\NormalTok{, }
               \AttributeTok{byrow =} \ConstantTok{FALSE}\NormalTok{) }
\NormalTok{mat1}
\end{Highlighting}
\end{Shaded}

\begin{verbatim}
     [,1] [,2] [,3] [,4]
[1,]    1    5    9   13
[2,]    2    6   10   14
[3,]    3    7   11   15
[4,]    4    8   12   16
\end{verbatim}

\begin{center}\rule{0.5\linewidth}{0.5pt}\end{center}

\subsection{Lists}\label{lists}

Unlike vectors and matrices, \textbf{lists} can contain elements of
different types. This makes them ideal for storing heterogeneous
information such as model outputs, metadata, or complex data structures.

\begin{Shaded}
\begin{Highlighting}[]
\CommentTok{\# Creating a list}
\NormalTok{my\_list }\OtherTok{\textless{}{-}} \FunctionTok{list}\NormalTok{(}\AttributeTok{name =} \StringTok{"John"}\NormalTok{, }\AttributeTok{age =} \DecValTok{25}\NormalTok{, }\AttributeTok{scores =} \FunctionTok{c}\NormalTok{(}\DecValTok{90}\NormalTok{, }\DecValTok{85}\NormalTok{, }\DecValTok{88}\NormalTok{))}
\FunctionTok{print}\NormalTok{(my\_list)}
\end{Highlighting}
\end{Shaded}

\begin{verbatim}
$name
[1] "John"

$age
[1] 25

$scores
[1] 90 85 88
\end{verbatim}

\begin{Shaded}
\begin{Highlighting}[]
\CommentTok{\# Accessing elements}
\NormalTok{my\_list}\SpecialCharTok{$}\NormalTok{name        }\CommentTok{\# by name}
\end{Highlighting}
\end{Shaded}

\begin{verbatim}
[1] "John"
\end{verbatim}

\begin{Shaded}
\begin{Highlighting}[]
\NormalTok{my\_list[[}\StringTok{"age"}\NormalTok{]]    }\CommentTok{\# by key}
\end{Highlighting}
\end{Shaded}

\begin{verbatim}
[1] 25
\end{verbatim}

\begin{Shaded}
\begin{Highlighting}[]
\NormalTok{my\_list[[}\DecValTok{3}\NormalTok{]][}\DecValTok{1}\NormalTok{]     }\CommentTok{\# first score}
\end{Highlighting}
\end{Shaded}

\begin{verbatim}
[1] 90
\end{verbatim}

Each element can itself be a vector, matrix, another list, or a data
frame.

\begin{center}\rule{0.5\linewidth}{0.5pt}\end{center}

\subsection{Summary}\label{summary}

\begin{itemize}
\tightlist
\item
  Data in R is stored as objects, and each has a \textbf{type} and
  \textbf{structure}.
\item
  Learn to inspect and convert types (\texttt{mode()}, \texttt{class()},
  \texttt{as.character()} etc.).
\item
  Vectors and matrices are efficient for numeric data, while lists and
  data frames offer more flexibility.
\item
  Naming and structuring your data properly will save you headaches down
  the line.
\item
  Factors help you manage categorical data and are essential for
  plotting and modelling. ````
\end{itemize}

\subsection{Practice Exercises}\label{practice-exercises}


\includegraphics{pics/Laptop_duck.webp}




\end{document}
